\section*{Resumo}

Este trabalho examina se o diagnóstico da literatura econômica sobre o impacto da inteligência artificial (IA) no mercado de trabalho mudou com o avanço da IA generativa. Conduz-se uma revisão sistemática, segundo as diretrizes PRISMA 2020, de estudos publicados entre 2013 e meados de 2026 e indexados na Web of Science e na Scopus. De 3.118 registros, 816 estudos foram incluídos e codificados, com apoio de modelos de linguagem na triagem e na extração, quanto ao tipo de evidência, à tecnologia em foco, aos mecanismos do arcabouço de tarefas de Acemoglu e Restrepo, ao sinal do efeito sobre o emprego e ao grupo de qualificação em risco. A fração de estudos que situam o risco na alta qualificação passa de 2,7\% entre os publicados até 2022 para 13,6\% entre os publicados de 2023 em diante, mas o risco para a baixa qualificação permanece o diagnóstico mais frequente. A mudança acompanha o objeto tecnológico, e não o calendário: o risco para a alta qualificação é a categoria modal entre os estudos sobre IA generativa (48,5\%) e é residual e estável entre os estudos sobre robótica e automação, usados como grupo de comparação. O resultado confirma a abordagem de tarefas, pois o risco segue o conteúdo das tarefas que cada tecnologia executa. Outras mudanças, como o horizonte mais curto e o sinal menos negativo, ocorrem também na literatura de robótica e refletem o amadurecimento empírico do campo. O sinal atribuído ao efeito sobre o emprego depende fortemente do tipo de evidência: estudos de exposição ocupacional tendem ao negativo, estudos de firma ao positivo, e modelos teóricos e revisões ao ambíguo. O novo diagnóstico apoia-se, por ora, em medidas de exposição e em revisões, não em efeitos observados. Discutem-se as limitações de cobertura da busca e do uso de modelos de linguagem como extratores.

\textbf{Palavras-chave:} inteligência artificial; IA generativa; mercado de trabalho; emprego; abordagem de tarefas; revisão sistemática.
\newpage
